% open-logic.tex
% compila o texto completo usando o estilo open-logic-dev

\documentclass[../include/open-logic-part]{subfiles}

\begin{document}

\clearpage

\begin{editorial}
Este arquivo carrega todo o conteúdo incluído no Open Logic Project.
Notas editoriais como esta, quando exibidas, indicam que o arquivo foi
compilado sem considerar como o material será apresentado. Se você
consegue ler esta nota, provavelmente \emph{não é aconselhável} ensinar
ou estudar a partir deste PDF.

O Open Logic Project oferece vários mecanismos para gerar um texto mais
adequado ao ensino ou ao estudo autônomo. Por padrão, por exemplo, o texto
trata todos os operadores lógicos como primitivos e desenvolve, nas provas,
todos os casos relativos a todos os operadores. Em geral, porém, é melhor
deixar alguns desses casos como exercícios. O Open Logic Project também é
um trabalho em andamento. Para estimular a colaboração e o aperfeiçoamento,
inclui-se material mesmo quando ele ainda está em forma de rascunho, não tem
exercícios etc. Um PDF preparado para um curso excluirá essas seções.

Para encontrar PDFs mais adequados ao ensino e ao estudo, consulte os
\href{http://builds.openlogicproject.org/}{cursos de exemplo disponíveis
  no site do OLP}. Para preparar sua própria versão, você pode começar pelo
\href{https://github.com/OpenLogicProject/OpenLogic/tree/master/courses/sample}{arquivo
  controlador de exemplo} ou consultar os fontes dos livros derivados para
ver exemplos mais elaborados e avançados.
\end{editorial}

\olimport[sets-functions-relations]{sets-functions-relations-complete}

\olimport[propositional-logic]{propositional-logic}

\olimport[first-order-logic]{first-order-logic}

\olimport[model-theory]{model-theory}

\olimport[computability]{computability}

\olimport[turing-machines]{turing-machines}

\olimport[incompleteness]{incompleteness}

\olimport[second-order-logic]{second-order-logic}

\olimport[lambda-calculus]{lambda-calculus}

\olimport[many-valued-logic]{many-valued-logic}

\olimport[normal-modal-logic]{normal-modal-logic}

\olimport[applied-modal-logic]{applied-modal-logic}

\olimport[intuitionistic-logic]{intuitionistic-logic}

\olimport[counterfactuals]{counterfactuals}

\olimport[set-theory]{set-theory}

\olimport[methods]{methods}

\olimport[history]{history}

\olimport[reference]{reference}

\end{document}

